% !TeX root = ../papers/nc-p3.tex
\section{一次元の点対応の次数と偏極}
\label{sec:point-graph-curves}
有限双対が標準Frobenius型で、四辺の左右射影ファイバーが空または一点、
一辺の像の次元が1とする。命題\ref{prop:surface-degree-bound}で全辺を共通の
スキーム $X$ に同定する。位置 $i$ の超平面束の引戻しを $L_i$ とする。
\begin{theorem}\label{thm:point-graph-curve-degree}
$X$ の一次元部分の総次数は、どちらの射影でも重複込みで高々7である。
任意の既約成分上で $\deg L_i$ は位置に依存しない。また $\operatorname{Pic}(X)$ で
\[
 L_i L_{i+2}^3=L_{i+1}^3 L_{i+3}.
\]
特に滑らかな楕円四次成分は存在すれば唯一で、その上で
$K=L_{i+2}L_i L_{i+1}^{-2}$ は $i$ に依存せず $K^{12}=\mathcal O$ である。
\end{theorem}
\begin{proof}
$B_j=\prod_{a=i+j}^{i+3} L_a$ と置き、$B_4=\mathcal O$ とする。
四文字のKoszul候補複体を点列上で評価すると、その各項は右端から順に
$B_j$ の $1,4,6,4,1$ 個の直和である。両端の微分は文字の点評価とその双対で階数1、
中央二微分は隣接関係の点評価と相補位置のその転置で階数3である。
複体条件と各ファイバーの完全性から局所分裂完全となる。
従って行列式の交代積は自明であり、表示した束の式を得る。
次数の三階差が零で、各成分の位置ごとの次数は二次以下の多項式である。
標準四移動は有限個の成分を置換し、各四位置の次数の集合は有限であるから、この多項式は定数である。

一辺の六関係が張る空間を $R$ とし、一般の超平面 $U\subset R$ を取る。
$U$ の五式の共通零点 $C\subset\mathbb P^3\times\mathbb P^3$ は完全交叉曲線である。
実際、元の基点集合の次元は1以下なので、一般の線型結合を順に選ぶと、
それ以前の交叉の次元2以上の成分を含むことはない。両射影は $C$ 上でも閉埋込みになる。
六式の片側評価が階数4の点で一式を除いても階数3以上である。階数3の点では、
階数を落とす $U$ は評価核（三次元）を含む超平面であり、これらは $\mathbb P^2$ をなす。
点自体は一次元以下の集合を動くので、両側についてこの閉集合の像は
$\mathbb P(R^*)$ を覆わない。残る $U$ の決定式軌跡で単元な三次小行列を使うと
次点の三座標比が正則に決まり、残りの式が射影像のスキームを定める。
随伴公式と交叉環より $\omega_C=\mathcal O_C(1,1)$、両次数10、算術種数11である。

$\mathcal O_X$ から零次元台の最大部分加群を除いた商の定める純曲線を $Y$ とする。
$C$ 内でそのイデアルの零化イデアルが定める残余曲線を $Z$ とする。
Gorenstein曲線上の双対性で $I_{Y/C}\otimes\omega_C=\omega_Z$ である。
実際、$\mathcal O_Y$ は極大Cohen--Macaulayで、その $\mathcal O_C$ 双対は
$\operatorname{ann} I_{Y/C}$、二重双対は元に戻る。イデアルの完全列を双対化すれば
残余商も純で、$\operatorname{ann} I_{Z/C}=I_{Y/C}$ を得る。
双対化層は各成分の一般点で商と同じ長さを持つため $\deg Y+\deg Z=10$ である。
第六式は $I_{Y/C}(1,1)$ の非零切断を与える。五式のKoszul解消を $(1,1)$ でひねると、
左の負の対角束は両射影空間のKünneth公式で次数5以下にコホモロジーを持たず、
$C$ を消す双一次式は $U$ だけだからである。Serre双対で $h^1(\mathcal O_Z)>0$ を得る。

射影三空間の純曲線で次数2以下なら $h^1(\mathcal O)=0$ である。
被約なら直線、平面二次曲線、または交わらない二直線である。
一本の直線 $T$ 上の二重構造では $I_T^2$ は一般点で零になり、純性により曲線上で零となる。
よって $I_{T/Z}$ は $T$ 上の階数1の無捩れ層で、$\mathcal O_T(-1)^2$ の商である。
これは次数 $\ge-1$ の直線束なので、構造層の完全列から同じ消滅が従う。
従って $\deg Z\ge3$ であり $\deg Y\le7$ である。

楕円四次成分は二本存在できず、標準四移動がその成分を保つ。
成分上で $D_i=L_{i+1}L_i^{-1}$ とすると $D_{i+1}=K D_i$。
四移動の引戻しが $\operatorname{Pic}^0$ に誘導する自己同型を $u$ として
$uK=K,\ uD_i=K^4D_i$（乗法記法）である。
$u=1$ なら $K^4=1$。他の場合 $u$ の位数は $2,3,4,6$ で、
$u-1$ の核は各々 $2,3,2,1$ で消える。これは各円分多項式を $1$ で評価すれば従う。
\end{proof}

\begin{corollary}\label{cor:point-graph-excess}
一辺のスキームが滑らかな連結曲線とこれに交わらない有限スキーム $F$ の直和であれば、
曲線の片側次数 $d$ と種数 $g$ に対して $\operatorname{length}F=18-4d+2g$ である。
従って楕円四次曲線なら残りの長さは4である。
\end{corollary}
\begin{proof}
六つの双一次式の最高Chern数は $\binom63=20$。
曲線に沿って六切断はイデアルを生成し、法束から $\mathcal O(1,1)^6$ への単射の余核は直線束である。
超過交叉公式で曲線の寄与はその次数となる。接束完全列より法束の次数は
$8d+2g-2$、余核の次数は $4d-2g+2$ である。
孤立部は滑らかな六次元空間で六式の完全交叉なので各局所長が寄与する。
\end{proof}
